\appendix
\chapter{Captures d'écran}

% Chaque capture est un fichier PNG du dossier captures/ ; tant qu'il manque, un encadré le signale.
\capture{catalogue}{Catalogue : liste paginée, filtre par catégorie et recherche}
\capture{fiche-produit}{Fiche produit (visiteur, client, produit en rupture)}
\capture{inscription-connexion}{Inscription et connexion, avec un message d'erreur}
\capture{panier}{Panier avec plusieurs lignes et total}
\capture{validation-commande}{Page de validation de la commande, puis message de confirmation}
\capture{stock-insuffisant}{Refus de commande pour stock insuffisant}
\capture{mes-commandes}{Historique « Mes commandes » et détail d'une commande avec le bouton d'annulation}
\capture{profil}{Profil du client}
\capture{tableau-de-bord}{Tableau de bord de l'administrateur}
\capture{admin-catalogue}{Back-office : catégories, formulaire produit, liste des produits (stock, produit inactif)}
\capture{admin-commandes}{Back-office : liste des commandes filtrée et changement de statut}
\capture{admin-utilisateurs}{Back-office : liste des utilisateurs}
\capture{page-erreur}{Page d'erreur (403 ou 404)}
\capture{mobile}{Affichage mobile (barre de navigation repliée)}
\capture{tables-mysql}{Tables générées dans MySQL}
\capture{github}{Historique Git et pull requests sur GitHub}

\chapter{Références}

\begin{itemize}
  \item Dépôt du projet : \url{https://github.com/unk-user/Boutique-JEE}
  \item Cahier des charges : \code{docs/cahier-des-charges.md}
  \item Diagrammes UML : \code{docs/uml/} (sources \code{.puml}, exports \code{.svg})
  \item Installation et déploiement : \code{README.md}
  \item Script du premier administrateur : \code{docs/sql/admin.sql}
\end{itemize}
